% =============================================================================
%  2-Desarrollo.tex — Estructura y contenido del taller (4 bloques)
% =============================================================================

\section{Estructura y contenido del taller}

El taller se organiza en cuatro bloques prácticos secuenciales, diseñados para
que cada paso construya sobre el anterior.  La duración total es de 2~horas.

% ─────────────────────────────────────────────────────────────────────────────
\subsection{Bloque 1 — Conexión e interoperabilidad (20 min)}

Configuración de la conexión nativa desde QGIS a los \textit{endpoints} de
Apimetro mediante protocolos HTTP(S).  Los participantes aprenderán a:

\begin{itemize}[nosep, leftmargin=1.5em]
  \item Agregar capas vectoriales remotas en QGIS utilizando la opción
    \textit{Protocol: HTTP(S)} y apuntando a las URLs de la API
    (\texttt{apimetro.dev}).
  \item Acceder de forma directa a geometrías GeoJSON estandarizadas de
    \textbf{líneas} (trazos de ruta tramo a tramo) y \textbf{estaciones}
    (nodos con jerarquía de transporte).
  \item Comprender la estructura del modelo de datos: entidades relacionales,
    campos de atributos y sistemas de referencia.
\end{itemize}

% ─────────────────────────────────────────────────────────────────────────────
\subsection{Bloque 2 — Visualización y contextualización (40 min)}

Extracción y superposición de las redes operativas multimodales sobre
cartografía base:

\begin{itemize}[nosep, leftmargin=1.5em]
  \item Carga de las capas de Metro, Metrobús, Mexibús, Cablebús, Trolebús,
    Tren Ligero, Tren Suburbano, RTP, Colectivo, Microbús y Autobús desde
    la API.
  \item Integración con cartografía base descargada mediante complementos de
    OpenStreetMap (QuickMapServices / QuickOSM) para evaluar el entorno
    construido.
  \item Simbología temática por sistema de transporte: colores diferenciados,
    categorización por jerarquía y etiquetado de nodos CETRAM.
\end{itemize}

% ─────────────────────────────────────────────────────────────────────────────
\subsection{Bloque 3 — Análisis espacial y cobertura (45 min)}

Aplicación de herramientas de geoprocesamiento para modelar accesibilidad
peatonal:

\begin{itemize}[nosep, leftmargin=1.5em]
  \item Generación de \textit{buffers} e isócronas de 800~m alrededor de cada
    estación para delimitar áreas de cobertura peatonal.
  \item Identificación cuantitativa de los vacíos de cobertura en las periferias
    metropolitanas ---especialmente en alcaldías como Milpa Alta, Tláhuac y
    municipios conurbados del Estado de México.
  \item Cálculo del factor de desviación ($DF$) en rutas específicas: cociente
    entre la distancia real de viaje en la red y la distancia euclidiana,
    siguiendo la metodología del \textit{VFTModel}
    \citep{garibaldi2026vft}.
\end{itemize}

% ─────────────────────────────────────────────────────────────────────────────
\subsection{Bloque 4 — Modelado y exportación (15 min)}

Métodos para almacenar y empaquetar las consultas dinámicas en formatos locales
estandarizados:

\begin{itemize}[nosep, leftmargin=1.5em]
  \item Exportación de las capas resultantes a GeoPackage (\texttt{.gpkg}) para
    asegurar la reproducibilidad de futuras investigaciones.
  \item Configuración de proyectos QGIS (\texttt{.qgz}) con las conexiones a
    la API preconfiguradas, listos para reutilizar.
  \item Discusión sobre flujos de trabajo reproducibles para proyectos de
    planeación urbana y tesis de posgrado.
\end{itemize}

% ─────────────────────────────────────────────────────────────────────────────
\subsection{Requisitos para los asistentes}

\begin{itemize}[nosep, leftmargin=1.5em]
  \item Computadora portátil con \textbf{QGIS 3.x} instalado.
  \item Conexión a internet (indispensable para las peticiones HTTP a la API).
  \item No se requieren conocimientos previos de programación.
\end{itemize}
